% Ad Familiares 5.8 (ad-familiares-05-08) --- English translation
% Translated by Claude (Anthropic) from the Latin (Perseus).
% Style guide: STYLE.md.

\ciceroLetterOpener{M. CICERO M. LICINIO P. F. CRASSO.}

\ciceroSection{1} How great my zeal has stood out for the defending or even augmenting of your standing, I do not doubt that all your people have written to you. For it was neither moderate, nor obscure, nor of the sort that could have been passed over in silence. For with both the consuls and with many consulars I fought in so great a strain as never before in any case; and I took on for myself a perpetual championship of all your honours, and a duty long owing to our old connection but interrupted by many shifts of the times, I have repaid in heaped measure.

\ciceroSection{2} Nor, by Hercules, has the will to cherish you or to do you honour ever been wanting in me; but certain plagues of men, those who grieve at another's praise, sometimes turned you against me, and sometimes changed me towards you. But a time has at last come, more wished for by me than hoped for: that, with your affairs most flourishing, both the memory of my goodwill and the faith of my friendship could be made manifest. For I have brought it to pass, not only that your whole household, but that the entire state, should know me to be most friendly to you. And so both your wife, the most outstanding of all women, and your sons, of distinguished dutifulness, virtue, and goodwill towards you, lean upon my counsels, admonitions, zeal, and actions; and the Senate and people of Rome understand that, in your absence, nothing is so prompt or so ready as my work, my care, my diligence, my authority in all things that concern you.

\ciceroSection{3} What things have been done and what are doing, I think you have learned by the letters of your own people. About myself I would have you think and would urgently wish you to persuade yourself thus: that I have not, by some sudden inclination or by chance, fallen upon embracing your eminence with my services, but that, from the moment I first touched the Forum, I have always looked to be as much as possible joined to you. From which time, indeed, I keep in memory that neither has my respect for you been wanting, nor your highest goodwill and generosity towards me. If anything has come up that was injured rather by suspicion than by fact, let those things, since they were both false and idle, be torn up out of all our memory and life. For you are such a man, and I desire to be such, that, since we have fallen on the same times of the commonwealth, I hope our union and friendship will be a praise to each of us.

\ciceroSection{4} Wherefore how much weight you think ought to be given to my judgement, you yourself will decide; and, as I hope, you will decide from a regard to my own standing. I indeed profess and pledge to you my exceptional and singular zeal in every kind of duty that bears upon your honour and your glory. In which, even if many will contend with me, yet --- with all the rest as judges, but especially with your two Crassi as judges --- I shall easily surpass them all. Both of whom I love uniquely; but, with equal goodwill toward Marcus, I am the more given over to Publius for this reason: that he --- though from his boyhood he has always done so --- yet at this time most of all both observes and loves me as a second parent.

\ciceroSection{5} I would have you think that this letter has the force of a treaty, not of a letter, and that what I promise and undertake to you I shall most religiously observe and most diligently perform. The defence of your standing in your absence, undertaken by me, I shall now persist in not only for the sake of our friendship but for the sake of my own consistency. Wherefore I have judged it enough, at this time, to write this to you: that, if I myself perceive anything that pertains either to your wish, your advantage, or your eminence, I shall do it of my own accord; but if I am ever advised of anything either by yourself or by your people, I shall bring it to pass that you understand that neither have you written, nor has any of your people referred anything to me, in vain. Wherefore I would have you write thus to me about all things great, small, and middling, as to a man most friendly to you; and instruct your people to use my service, counsel, authority, and goodwill in all your business, public, private, forensic, domestic, and in the business of your friends, guests, and clients, in such a way that, so far as it can be done, the longing for your presence may be lessened by my labour.
